\documentclass[../../../include/open-logic-section]{subfiles}

\begin{document}

\olfileid{sth}{story}{extensionality}
\olsection{外延性}

首先要指出的是，集合是由其元素唯一确定的。更精确地说：

\begin{axiom}[外延性公理]
如果集合 $A$ 和 $B$ 具有相同的元素，则 $A$ 和 $B$ 是同一个集合。
\[
  \lforall[A][\lforall[B][(\lforall[x][(x \in A \liff x \in B)] \lif
  \eq[A][B])]]
\]
\end{axiom}

我们在 \olref[sfr][][]{part} 中始终预设了这一点。但这点值得重申。外延性公理表达了集合由其元素决定这一基本思想。（因此，集合或许可以与\emph{概念}形成对比：完全相同的对象可以归属于许多不同的概念。）

为什么要采纳这条原则？可以说，任何对外延性的否定，都无异于决定放弃任何哪怕能被称为\emph{集合论}的东西。集合论无非是关于外延汇集的理论。

然而，在 \olref[sth][][]{part} 中，真正的挑战在于确立一些原则，用以告诉我们\emph{哪些集合存在}。而结果表明，对这个问题唯一真正“明显”的答案，是可以证明为错误的。

\end{document}